\chapter{Submanifold \texorpdfstring{$\R^n$}{Rn}}
\label{ch:b3:submanifolds}

Bola, torus, grup rotasi: habitat alami
geometri dan mekanika bukanlah ruang vektor melainkan \emph{himpunan
melengkung yang tampak datar bila didekati}. Bab ini memberi kalimat itu
makna yang persis --- yakni submanifold $\R^n$ --- beserta
kalkulus untuk bekerja padanya. Landasannya adalah teorema fungsi
invers, yang di sini dibuktikan lewat titik tetap Banach;
sedangkan selebihnya adalah penggantian koordinat: yakni keempat pemerian
submanifold yang setara (pelurusan lokal, himpunan aras,
grafik, parametrisasi), \omterm{def:b3:submanifolds:tangent}{ruang singgung}, dan pengoptimalan
terkendala lewat \omterm{thm:b3:submanifolds:lagrange}{pengali Lagrange} --- yang, sebagai peragaan
perpisahan, membuktikan kembali teorema spektral bagi matriks setangkup
dalam tiga baris geometri. Sedangkan soal akhir pekannya
membangun grup rotasi $SO(3)$ beserta selimut ganda kuaternioniknya:
yakni aljabar ($Q_8$ \cref{ch:b3:groups} yang telah dewasa) yang
berjumpa geometri.

\section{Teorema fungsi invers}

\begin{theorem}[Teorema fungsi invers]
\label{thm:b3:submanifolds:ift}
Misalkan $U \subseteq \R^n$ terbuka, $f \colon U \to \R^n$ berkelas
$\mathcal C^1$, dan $a \in U$ dengan $Df(a)$ yang terbalikkan.
Maka ada \omterm{def:b3:topology:topology}{himpunan terbuka} $V \ni a$ dan $W \ni f(a)$ sedemikian sehingga
$f \colon V \to W$ merupakan bijeksi dengan balikan $\mathcal C^1$,
dan\index{teorema fungsi invers}
\[
D(f^{-1})(y) = \bigl(Df(f^{-1}(y))\bigr)^{-1}
\qquad (y \in W).
\]
Dan jika $f$ bersifat $\mathcal C^k$, maka demikian pula $f^{-1}$.
\end{theorem}

\begin{proof}
Normalkanlah: dengan mengganti $f$ oleh $x \mapsto Df(a)^{-1}\bigl(f(a +
x) - f(a)\bigr)$, kita boleh mengandaikan $a = 0$, $f(0) = 0$, dan $Df(0) =
I$ (sebab pernyataan umumnya menyusul lewat pengomposisian dengan
bijeksi afinnya). Tulislah $f(x) = x + g(x)$: maka $Dg(0) = 0$, dan
lewat \omterm{def:b3:topology:continuity}{kekontinuan} $Dg$ pilihlah $r > 0$ dengan
$\vertiii{Dg(x)} \leq \frac12$ pada $\bar B(0, r)$; sehingga ketaksamaan
nilai rata-ratanya memberikan $\norm{g(x) - g(x')} \leq \frac12
\norm{x - x'}$ di sana.

\emph{Kebijektifannya ke sebuah \omterm{def:b3:topology:topology}{persekitaran}.} Untuk $y \in B(0,
\frac r2)$, menyelesaikan $f(x) = y$ berarti mencari titik tetap
$\Phi_y(x) = y - g(x)$; sedangkan $\Phi_y$ memetakan $\bar B(0, r)$ ke
dirinya (sebab $\norm{\Phi_y(x)} \leq \norm y + \frac12\norm x \leq
r$) dan bersifat Lipschitz-$\frac12$: sehingga Banach
(\cref{thm:b3:complete:banach}) memberikan satu-satunya solusi
$x = \varphi(y) \in \bar B(0,r)$. Lebih lanjut $f$ bersifat injektif
pada $\bar B(0,r)$:
\[
\norm{f(x) - f(x')} \geq \norm{x - x'} - \norm{g(x) - g(x')}
\geq \tfrac12\norm{x - x'} .
\tag{$*$}
\]
Lalu tetapkan $W = B(0, \frac r2)$ dan $V = f^{-1}(W)\cap B(0, r)$:
yang terbuka (lewat \omterm{def:b3:topology:continuity}{kekontinuannya}), dengan $f \colon V \to W$ yang bijektif.

\emph{\omterm{def:b3:topology:continuity}{Kekontinuan} dan keterturunan balikannya.} Ketaksamaan $(*)$
mengatakan bahwa $\varphi = f^{-1}$ bersifat Lipschitz-$2$. Tetapkan $y_0 = f(x_0)
\in W$; lalu keterbalikan $A = Df(x_0)$ (sebab jaraknya ke $I$
adalah $\leq \frac12$: menurut Neumann,
\cref{prop:b3:banach:neumann}) dan keterturunan $f$
memberikan, untuk $y = f(x)$ di dekat $y_0$:
\[
\varphi(y) - \varphi(y_0) - A^{-1}(y - y_0)
= -A^{-1}\bigl(f(x) - f(x_0) - A(x - x_0)\bigr)
= -A^{-1}\,o(\norm{x - x_0}) = o(\norm{y - y_0}),
\]
dengan memakai $(*)$ untuk mengubah $\norm{x - x_0} \leq 2\norm{y - y_0}$:
sehingga $\varphi$ dapat diturunkan di $y_0$ dengan diferensial
balikannya. Sedangkan \omterm{def:b3:topology:continuity}{kekontinuan} $y \mapsto D\varphi(y) =
Df(\varphi(y))^{-1}$: lewat komposisi \omterm{def:b3:topology:continuity}{pemetaan kontinu}
(sebab pembalikannya \omterm{def:b3:topology:continuity}{kontinu},
\cref{prop:b3:banach:neumann}): jadi $\varphi \in \mathcal C^1$;
dan mengangkat rumus yang sama memberikan $\mathcal C^k$.
\end{proof}

\begin{theorem}[Teorema fungsi implisit]
\label{thm:b3:submanifolds:implicit}
Misalkan $F \colon U \subseteq \R^p\times\R^q \to \R^q$ bersifat
$\mathcal C^1$ di dekat $(a, b)$, $F(a,b) = 0$, dan andaikan
diferensial parsial $D_yF(a,b) \in \mathcal L(\R^q)$ bersifat
terbalikkan. Maka ada \omterm{def:b3:topology:topology}{persekitaran} $A \ni a$ dan $B \ni b$
beserta sebuah pemetaan $\mathcal C^1$ $\psi \colon A \to B$ dengan
\[
\bigl\{(x, y)\in A\times B : F(x,y) = 0\bigr\}
= \{(x, \psi(x)) : x \in A\},
\]
dan $D\psi(x) = -D_yF(x, \psi(x))^{-1}\,D_xF(x,
\psi(x))$.\index{teorema fungsi implisit}
\end{theorem}

\begin{proof}
Terapkan \cref{thm:b3:submanifolds:ift} pada $\Theta(x, y) = (x,
F(x,y))$: sebab diferensialnya di $(a,b)$, yang segitiga blok dengan
blok diagonal $I$ dan $D_yF$ yang terbalikkan, bersifat terbalikkan. Lalu
balikan lokalnya berbentuk $\Theta^{-1}(x, z) = (x, h(x,
z))$; lalu tetapkan $\psi(x) = h(x, 0)$: maka $F(x, y) = 0$ jika dan hanya jika
$\Theta(x,y) = (x, 0)$ jika dan hanya jika $y = \psi(x)$, secara lokal. Sedangkan
rumusnya: turunkanlah $F(x, \psi(x)) = 0$ lewat aturan
rantainya.
\end{proof}

\section{Submanifold: empat definisi}

\begin{theorem}[Pencirian yang setara]
\label{thm:b3:submanifolds:characterizations}
Misalkan $M \subseteq \R^n$, $d \in \{0, \dots, n\}$, dan $k \geq
1$. Pernyataan berikut setara, untuk setiap titik $a \in M$
(dan $M$ disebut \emph{submanifold berdimensi $d$ berkelas
$\mathcal C^k$}\index{submanifold} jika keduanya berlaku di setiap $a
\in M$):
\begin{enumerate}
  \item (\emph{Pelurusan}) Ada sebuah difeomorfisma $\mathcal C^k$
        $\Phi$ dari sebuah $\Omega \ni a$ yang terbuka
        ke sebuah $\Omega' \subseteq \R^n$ yang terbuka dengan
        \[
        \Phi(M\cap\Omega) = \Omega' \cap
        \bigl(\R^d\times\{0\}\bigr).
        \]
  \item (\emph{Himpunan aras}) Ada sebuah \emph{submersi}
        $\mathcal C^k$ $F \colon \Omega \to \R^{n-d}$
        (yakni $DF(x)$ yang surjektif) pada sebuah $\Omega \ni a$ yang terbuka
        dengan $M\cap\Omega = F^{-1}(0)$.
  \item (\emph{Grafik}) Sampai permutasi koordinatnya, $M$ secara
        lokal merupakan grafik sebuah pemetaan $\mathcal C^k$ $\psi
        \colon A \subseteq \R^d \to \R^{n-d}$.
  \item (\emph{Parametrisasi}) Ada sebuah \emph{imersi}
        $\mathcal C^k$ $\varphi \colon A \subseteq \R^d \to
        \R^n$ (dengan $D\varphi(u)$ yang injektif) dan $A$ yang terbuka,
        sedemikian sehingga $\varphi$ merupakan \omterm{def:b3:topology:continuity}{homeomorfisma} dari $A$ ke $M \cap
        \Omega$ untuk suatu $\Omega \ni a$ yang terbuka.
\end{enumerate}
\end{theorem}

\begin{proof}
(1)$\Rightarrow$(2): ambillah $F = (\Phi_{d+1}, \dots, \Phi_n)$ (yakni koordinat
terakhir $\Phi$): sebuah submersi (sebab $D\Phi$ terbalikkan).
(2)$\Rightarrow$(3): karena $DF(a)$ surjektif, suatu minor $q \times q$
Jacobinya bersifat terbalikkan (dengan $q = n - d$); lalu setelah
koordinatnya dipermutasikan, $D_yF(a)$ terbalikkan, dan teorema fungsi implisitnya (\cref{thm:b3:submanifolds:implicit})
menyatakan $M$ secara lokal sebagai grafik $y = \psi(x)$.
(3)$\Rightarrow$(4): ambillah $\varphi(x) = (x, \psi(x))$: yakni sebuah
imersi (sebab diferensialnya $\bigl(\begin{smallmatrix}I\\
D\psi\end{smallmatrix}\bigr)$ injektif), sekaligus \omterm{def:b3:topology:continuity}{homeomorfisma} ke
grafiknya (dengan balikannya: proyeksinya, yang \omterm{def:b3:topology:continuity}{kontinu}).
(4)$\Rightarrow$(1): misalkan $\varphi(u_0) = a$; lalu lengkapkanlah
$\operatorname{im}D\varphi(u_0)$ oleh sebuah pelengkap $E$
(dengan $\dim E = n - d$) lalu definisikan $\Theta(u, v) = \varphi(u) + v$
pada $A\times E$: maka $D\Theta(u_0, 0)$ bersifat bijektif (sebab petanya
memuat $\operatorname{im}D\varphi(u_0)$ dan $E$), sehingga
$\Theta$ merupakan difeomorfisma lokal
(\cref{thm:b3:submanifolds:ift}); dan balikannya $\Phi$
meluruskan: sebab di dekat $a$, titik $M$ tepat berupa
$\varphi(u) = \Theta(u, 0)$ --- dan untuk ini, hipotesis \omterm{def:b3:topology:continuity}{homeomorfisma}
pada (4) menjamin bahwa $M\cap\Omega$, untuk $\Omega$ yang kecil,
tak memuat lembar lain (sebab $\varphi(u') + v = m \in
M$ yang dekat dengan $a$ dengan $v \neq 0$ yang kecil harus dikecualikan:
karena $m = \varphi(u'')$ untuk suatu $u''$ di dekat $u_0$ menurut
sifat \omterm{def:b3:topology:continuity}{homeomorfismanya}, dan keinjektifan lokal $\Theta$
memaksa $v = 0$). Lalu $\Phi(M\cap\Omega) = (A\times\{0\})
\cap\Phi(\Omega)$ sampai penyusutannya.
\end{proof}

\begin{example}\label{ex:b3:submanifolds:examples}
Bola $S^{n-1} = \{\norm x^2 = 1\}$: yakni himpunan aras
submersi $F(x) = \norm x_2^2 - 1$ pada
$\R^n\setminus\{0\}$ (sebab $DF(x) = 2x^{\mathsf T} \ne 0$):
jadi submanifold $\mathcal C^\infty$ berdimensi $n - 1$. Sedangkan
torus di $\R^3$: yakni himpunan aras $\bigl(\sqrt{x^2 + y^2} -
R\bigr)^2 + z^2 - r^2$ (dengan $0 < r < R$). Sementara kerucut $\{x^2 + y^2 =
z^2\}$ \emph{bukan} submanifold di $0$
(\cref{exo:b3:submanifolds:1}). Untuk grup matriksnya: $SL_n$ dan
$O_n$ merupakan submanifold $M_n(\R)$
(\cref{exo:b3:submanifolds:5,exo:b3:submanifolds:6}) --- yakni
titik awal teori Lie.
\end{example}

\section{Ruang singgung}

\begin{definition}\label{def:b3:submanifolds:tangent}
Misalkan $M$ sebuah submanifold-$d$ dan $a \in M$. Adapun \emph{ruang
singgung}\index{ruang singgung} $T_aM$ adalah himpunan vektor
kecepatan $\gamma'(0)$ milik kurva $\mathcal C^1$ $\gamma \colon
\intoo{-\varepsilon}\varepsilon \to M$ dengan $\gamma(0) = a$.
\end{definition}

\begin{proposition}\label{prop:b3:submanifolds:tangent}
Ruang $T_aM$ merupakan subruang berdimensi $d$ dari $\R^n$, dan:
\begin{enumerate}
  \item jika $M = F^{-1}(0)$ secara lokal dengan $F$ sebuah submersi:
        maka $T_aM = \ker DF(a)$;
  \item jika $M$ terparametrikan oleh imersi $\varphi$
        (dengan $\varphi(u_0) = a$): maka $T_aM =
        \operatorname{im}D\varphi(u_0)$.
\end{enumerate}
\end{proposition}

\begin{proof}
Kurva di $M$ memenuhi $F(\gamma(t)) = 0$; sehingga aturan rantai di
$0$ memberikan $DF(a)\gamma'(0) = 0$: jadi $T_aM \subseteq \ker DF(a)$.
Sebaliknya, pelurusannya (\cref{thm:b3:submanifolds:%
characterizations}(1)) mengangkut garis
$\R^d\times\{0\}$ menjadi kurva $M$: sehingga setiap vektor sebuah
subruang berdimensi $d$ terwujud; lalu membandingkan dimensinya
(sebab $\dim\ker DF(a) = n - (n - d) = d$) memaksa kesamaannya pada (1),
dan argumen pengangkutan yang sama memberikan (2) (yakni $D\varphi(u_0)$
yang diterapkan pada garis lurus di $A$; lewat dimensinya lagi).
\end{proof}

\begin{theorem}[Pengali Lagrange]
\label{thm:b3:submanifolds:lagrange}
Misalkan $M = F^{-1}(0)$ dengan $F = (F_1, \dots, F_q) \colon \Omega
\to \R^q$ sebuah submersi $\mathcal C^1$, dan misalkan $f \colon
\Omega \to \R$ bersifat $\mathcal C^1$. Jika pembatasan
$f\restriction_M$ berekstremum lokal di $a \in M$, maka
ada real yang tunggal $\lambda_1, \dots, \lambda_q$
(yakni \emph{\omterm{thm:b3:submanifolds:lagrange}{pengali Lagrange}}) dengan\index{pengali Lagrange}
\[
\nabla f(a) = \lambda_1\nabla F_1(a) + \dots +
\lambda_q\nabla F_q(a) .
\]
\end{theorem}

\begin{proof}
Untuk setiap kurva $\gamma$ di $M$ yang melalui $a$: $t \mapsto
f(\gamma(t))$ berekstremum lokal di $0$, sehingga $0 =
\frac{\dd}{\dd t}f(\gamma(t))\big|_0 = \langle\nabla f(a),
\gamma'(0)\rangle$: jadi $\nabla f(a) \perp T_aM = \ker DF(a)$
(\cref{prop:b3:submanifolds:tangent}). Kini $\ker DF(a)^\perp
= \operatorname{im}DF(a)^{\mathsf T}$ (yakni identitas rank dan
keortogonalan Tahun ke-2, atau \cref{exo:b3:hilbert:8} dalam dimensi
berhingga: $(\ker T)^\perp = \operatorname{im}T^*$), yang
direntang gradiennya $\nabla F_i(a)$ --- yang bebas, sebab
$DF(a)$ surjektif: sehingga pengalinya ada dan tunggal.
\end{proof}

\begin{example}[Teorema spektralnya, secara geometri]
\label{ex:b3:submanifolds:rayleigh}
Misalkan $A$ sebuah matriks setangkup real $n\times n$ lalu maksimumkan
$f(x) = \langle Ax, x\rangle$ pada bola $S^{n-1}$
(yang \omterm{def:b3:topology:compact}{kompak}: sehingga maksimumnya tercapai, di suatu $v_1$). Lalu Lagrange
dengan $F(x) = \norm x^2 - 1$: $\nabla f = 2Ax$ dan $\nabla F =
2x$ memberikan $Av_1 = \lambda_1v_1$ --- yakni sebuah vektor eigen, dengan
$\lambda_1 = \max_{S^{n-1}}\langle Ax, x\rangle$. Lalu batasilah
$A$ ke $v_1^\perp$ (yang invarian: sebab $\langle Av, v_1\rangle =
\langle v, Av_1\rangle = \lambda_1\langle v, v_1\rangle = 0$)
lalu iterasikan: diperolehlah basis ortonormal berisi vektor eigen. Jadi
teorema spektral Tahun ke-2, yang dibuktikan kembali lewat pengoptimalan murni
--- dan bayangan berdimensi tak berhingga argumen yang sama
membuktikan \cref{lem:b3:spectral:existence}.
\end{example}

\begin{method}\label{met:b3:submanifolds:method}
Untuk membuktikan sebuah himpunan merupakan submanifold: tampilkanlah ia secara lokal sebagai
$F^{-1}(0)$ dengan $DF$ yang surjektif \emph{pada himpunannya} (yakni jalur yang
paling lazim), atau sebagai sebuah grafik. Untuk menghitung dimensi dan
\omterm{def:b3:submanifolds:tangent}{ruang singgungnya}: $d = n - q$ dan $T_a = \ker DF(a)$. Untuk
mengoptimalkan padanya: pakailah Lagrange --- dan periksalah selalu kekompakannya (atau
kekoersifannya) lebih dahulu, agar ada ekstremum yang teoremanya
dapat berlaku padanya, lalu ingatlah bahwa persamaan pengalinya hanya
\emph{perlu}: kumpulkanlah semua titik kritisnya, lalu bandingkan
nilainya. Sedangkan untuk grup matriksnya, turunkanlah kurvanya di
identitasnya untuk mengenali \omterm{def:b3:submanifolds:tangent}{ruang singgungnya}.
\end{method}

\section{Latihan}

\begin{exercise}[$\star$]\label{exo:b3:submanifolds:1}
(a) Periksalah bahwa berikut ini merupakan submanifold $\mathcal C^\infty$
lalu berikan dimensinya: $S^{n-1}$;
hiperboloid $\{x^2 + y^2 - z^2 = 1\}$; dan torus pada
\cref{ex:b3:submanifolds:examples}.
(b) Tunjukkan bahwa kerucut $C = \{x^2 + y^2 = z^2\} \subseteq
\R^3$ bukan submanifold-$2$ di $0$: tentukanlah banyaknya
\omterm{def:b3:topology:components}{komponen terhubung}
$\bigl(C\setminus\{0\}\bigr)\cap B(0,\varepsilon)$, lalu
bandingkan dengan banyaknya yang akan dipaksa sebuah pelurusan
(\cref{thm:b3:submanifolds:characterizations}(1))
bagi sebuah bidang dikurangi sebuah titik.
\end{exercise}

\begin{exercise}[$\star$]\label{exo:b3:submanifolds:2}
Hitunglah \omterm{def:b3:submanifolds:tangent}{ruang singgungnya}:
(a) $T_aS^{n-1}$ untuk sembarang $a$ (jawabannya: $a^\perp$);
(b) bidang singgung torus pada
\cref{ex:b3:submanifolds:examples} di sembarang titik
khatulistiwa luarnya $\{z = 0,\ x^2 + y^2 = (R + r)^2\}$;
(c) garis singgung heliks $\varphi(t) = (\cos t,
\sin t, t)$ di $\varphi(t_0)$, sambil memeriksa
\cref{prop:b3:submanifolds:tangent}(2).
\end{exercise}

\begin{exercise}[$\star\star$]\label{exo:b3:submanifolds:3}
Misalkan $f(x, y) = (x^2 - y^2,\ 2xy)$ (yakni $z \mapsto z^2$).
(a) Di titik mana saja \cref{thm:b3:submanifolds:ift} berlaku?
(b) Tunjukkan bahwa $f$ terbalikkan secara lokal tetapi tak secara global pada
$\R^2\setminus\{0\}$, lalu tampilkan secara eksplisit kedua balikan
lokalnya yang terdefinisi pada sebuah \omterm{def:b3:topology:topology}{persekitaran} $(1, 0)$ (yakni kedua
cabang akar kuadratnya).
(c) Lakukan pembahasan yang sama bagi pemetaan \omterm{ex:b3:product:polar}{koordinat kutub} $(r,
\theta) \mapsto (r\cos\theta, r\sin\theta)$.
\end{exercise}

\begin{exercise}[$\star\star$]\label{exo:b3:submanifolds:4}
(Folium) Misalkan $F(x,y) = x^3 + y^3 - 3xy$ dan $\mathcal C =
F^{-1}(0)$.
(a) Tunjukkan bahwa di dekat setiap titik $\mathcal C$ selain
titik asalnya, $\mathcal C$ merupakan submanifold-$1$ yang secara lokal berupa grafik terhadap
$x$ atau terhadap $y$ (yang mana, di mana?).
(b) Hitunglah garis singgungnya di $(\frac32, \frac32)$.
(c) Apa yang terjadi di titik asalnya? (Sebab dua cabangnya bersilangan:
tampilkanlah dua kurva $\mathcal C^1$ di $\mathcal C$ yang melalui $0$
dengan kecepatan yang bebas, lalu simpulkan bahwa tak ada
pelurusan yang ada.)
\end{exercise}

\begin{exercise}[$\star\star$]\label{exo:b3:submanifolds:5}
Misalkan $F(M) = M^{\mathsf T}M$ dari $M_n(\R)$ ke ruang
$S_n$ berisi matriks setangkup.
(a) Tunjukkan $DF(M)(H) = M^{\mathsf T}H + H^{\mathsf T}M$ dan
bahwa $DF(M)$ bersifat surjektif ke $S_n$ di setiap $M \in O_n$
\emph{(diberikan $S \in S_n$, cobalah $H = \frac12MS$)}.
(b) Simpulkan bahwa $O_n = F^{-1}(I)$ merupakan submanifold
$\mathcal C^\infty$ yang \omterm{def:b3:topology:compact}{kompak} berdimensi
$\frac{n(n-1)}2$, dengan $T_IO_n = \{H : H^{\mathsf T} = -H\}$,
yakni matriks antisetangkupnya.
(c) Tunjukkan bahwa $\eu^{tH} \in O_n$ untuk setiap $H$ yang antisetangkup:
sehingga arah singgungnya \omterm{def:b3:lebesgue:l1}{terintegralkan} menjadi kurva di grupnya.
\end{exercise}

\begin{exercise}[$\star\star$]\label{exo:b3:submanifolds:6}
(a) Tunjukkan bahwa $\det \colon M_n(\R) \to \R$ berdiferensial
$D\det(M)(H) = \operatorname{tr}\bigl(\operatorname{com}(M)
^{\mathsf T}H\bigr)$, yang tak nol di setiap $M \in SL_n$.
(b) Simpulkan bahwa $SL_n(\R)$ merupakan submanifold berdimensi
$n^2 - 1$ dengan $T_ISL_n = \{H : \operatorname{tr}H = 0\}$.
(c) Apakah $GL_n(\R)$ sebuah submanifold? Berdimensi berapa?
\end{exercise}

\begin{exercise}[$\star\star$]\label{exo:b3:submanifolds:7}
Lewat \omterm{thm:b3:submanifolds:lagrange}{pengali Lagrange}:
(a) carilah ekstremum $f(x,y) = xy$ pada lingkaran $x^2 +
y^2 = 1$;
(b) tunjukkan bahwa di antara semua vektor peluang $(p_1, \dots,
p_n)$ (yang positif dan berjumlah $1$), entropi $-\sum
p_i\ln p_i$ dimaksimumkan tepat pada distribusi seragamnya;
(c) carilah titik elips $\{x^2/4 + y^2 = 1\}$ yang
terdekat ke $(1, 0)$, lalu periksalah persamaan pengalinya
secara geometri (lewat kesejajaran normalnya).
\end{exercise}

\begin{exercise}[$\star\star\star$]\label{exo:b3:submanifolds:8}
Tuliskanlah \cref{ex:b3:submanifolds:rayleigh} selengkapnya: buktikan
lewat induksi bahwa matriks setangkup real memiliki
basis ortonormal berisi vektor eigen, dengan
$\lambda_1 \geq \dots \geq \lambda_n$ sebagai maksimum
terkendala berturutan hasil bagi Rayleighnya. Lalu turunkan
rumus Courant--Fischer \cref{exo:b3:spectral:8} dalam
dimensi berhingga langsung dari konstruksi ini.
\end{exercise}

\begin{exercise}[$\star\star\star$]\label{exo:b3:submanifolds:9}
(Ketaksamaan Hadamard) Untuk $M \in GL_n(\R)$ yang berkolom
$c_1, \dots, c_n$:
\[
\abs{\det M} \;\leq\; \prod_{i=1}^n\norm{c_i}_2 ,
\]
dengan kesamaannya jika dan hanya jika kolomnya ortogonal. \emph{(Susutkanlah
ke kolom bernorma $1$ lewat penskalaan; lalu maksimumkan $\det$ pada
hasil kali bola yang \omterm{def:b3:topology:compact}{kompak} $(S^{n-1})^n$; lalu pada sebuah pemberi maksimum,
Lagrange pada setiap kolomnya secara terpisah memberikan $\nabla_{c_i}\det =
\lambda_ic_i$, sedangkan $\nabla_{c_i}\det$ adalah kolom ke-$i$
$\operatorname{com}(M)$: lalu turunkan bahwa $M^{\mathsf T}M$ diagonal,
sehingga $= I$, sehingga $\det = \pm1$.)} Bacaan geometrinya: bahwa
volume sebuah paralelepipedum paling banyak hasil kali panjang
rusuknya.
\end{exercise}

\begin{exercise}[$\star\star$]\label{exo:b3:submanifolds:10}
Di dekat titik yang mana lingkaran $S^1$ berupa grafik $y =
\psi(x)$? Dan grafik $x = \chi(y)$? Periksalah pencirian
grafiknya (\cref{thm:b3:submanifolds:%
characterizations}(3)) secara eksplisit di $(1, 0)$, lalu terangkan
dalam satu kalimat mengapa \emph{suatu} permutasi koordinat
selalu mencukupi tetapi tak ada satu pun yang selalu berhasil.
\end{exercise}

\begin{exercise}[$\star\star$]\label{exo:b3:submanifolds:11}
(Grup ortogonal sebagai submanifold, secara kuantitatif)
(a) Tunjukkan bahwa $O_n = \{M : M^{\mathsf T}M = I\}$ bersifat \omterm{def:b3:topology:compact}{kompak}:
yakni terbatas (sebab setiap kolomnya vektor satuan, sehingga $\norm M \leq
\sqrt n$ untuk norma matriks Euclidnya) dan tertutup.
(b) Tunjukkan bahwa \omterm{def:b3:submanifolds:tangent}{ruang singgungnya} di $I$ adalah ruang berisi matriks
antisetangkup, yang berdimensi $\frac{n(n-1)}2$, dan
di sebuah $A \in O_n$ yang umum: $T_AO_n = \{AK : K^{\mathsf T} =
-K\}$.
(c) Turunkan bahwa pemetaan $t \mapsto A\exp(tK)$ merupakan, untuk setiap
$K$ yang antisetangkup, sebuah kurva di $O_n$ yang melalui $A$ dengan
kecepatan $AK$ \emph{(periksalah $\exp(tK) \in O_n$ dengan memakai
$\exp(X)^{\mathsf T} = \exp(X^{\mathsf T})$ dan
$\exp(-X)\exp(X) = I$)}: sehingga setiap vektor singgungnya terwujud oleh
sebuah kurva yang eksplisit, tanpa memerlukan teorema fungsi implisit.
\end{exercise}

\begin{exercise}[$\star\star$]\label{exo:b3:submanifolds:12}
(Titik kritis jaraknya) Misalkan $M \subseteq \R^n$ sebuah
submanifold dan $p \notin M$. Tunjukkan bahwa jika $x_0 \in M$
meminimumkan jaraknya ke $p$ (dan titik semacam itu ada bila $M$
tertutup dan tak kosong --- mengapa?), maka
\[
p - x_0 \;\perp\; T_{x_0}M
\]
\emph{(turunkanlah $t \mapsto \norm{\gamma(t) - p}^2$
sepanjang kurva di $M$)}. Lalu turunkan: bahwa titik terdekat pada
sebuah bola terletak pada sinar yang melalui pusatnya; lalu pakailah
syaratnya untuk menghitung jarak dari $p = (2, 0)$ ke
parabola $y = x^2$ (susutkanlah ke sebuah kubik lalu selesaikan
secara numerik sampai tiga angka).
\end{exercise}

\section{Soal: \texorpdfstring{$SO(3)$}{grup rotasi} dan
kuaternion}

\begin{problem}[{Soal akhir pekan --- rotasi, grup
$S^3$, dan selimut gandanya}]\label{pb:b3:submanifolds:1}
Kuaternion $\mathbb H = \{t + x\mathrm i + y\mathrm j +
z\mathrm k\}$ --- yakni aljabar yang grup satuannya memuat
$Q_8$ \cref{pb:b3:groups:1} --- memparametrikan
rotasi tiga dimensi dua kali lipat: sebab pemetaan ``konjugatkan
dengan sebuah kuaternion satuan'' merupakan morfisma surjektif $S^3 \to
SO(3)$ berkernel $\{\pm1\}$. Kita membangun segalanya.
Ingat kembali atau definisikan: perkaliannya bilinear-$\R$ dengan
$\mathrm i^2 = \mathrm j^2 = \mathrm k^2 =
\mathrm{ijk} = -1$; sedangkan konjugat $q = t + x\mathrm i +
y\mathrm j + z\mathrm k$ adalah $\bar q = t - x\mathrm i -
y\mathrm j - z\mathrm k$; dan $N(q) = q\bar q = t^2 + x^2 + y^2
+ z^2$.

\medskip
\noindent\textbf{Bagian I --- Aljabar $\mathbb H$ dan
grup $S^3$.}
\begin{enumerate}
  \item Periksalah bahwa $\mathbb H$ merupakan aljabar-$\R$ yang
        asosiatif berpusat $\R$, bahwa $\overline{pq} =
        \bar q\,\bar p$, dan bahwa $N(pq) = N(p)N(q)$
        \emph{(salah satu jalur yang bersih: sajikanlah $q$ sebagai
        matriks kompleks $2\times2$
        $\bigl(\begin{smallmatrix}\alpha & \beta\\
        -\bar\beta & \bar\alpha\end{smallmatrix}\bigr)$ dengan $q =
        \alpha + \beta\mathrm j$, lalu pakailah $\det$)}.
  \item Turunkan bahwa setiap $q \neq 0$ bersifat terbalikkan (sebab $q^{-1} =
        \bar q/N(q)$): sehingga $\mathbb H$ merupakan lapangan (yang tak
        komutatif), dan $S^3 = \{N(q) = 1\}$ merupakan grup ---
        sekaligus submanifold-$3$ $\R^4$ yang \omterm{def:b3:topology:compact}{kompak}
        (\cref{ex:b3:submanifolds:examples}).
\end{enumerate}

\medskip
\noindent\textbf{Bagian II --- Morfisma rotasinya.} Kenalilah
$\R^3$ dengan \emph{kuaternion murni} $P = \{x\mathrm i +
y\mathrm j + z\mathrm k\}$, lalu untuk $q \in S^3$ definisikan
$\rho_q(v) = q\,v\,\bar q$.
\begin{enumerate}[resume]
  \item Tunjukkan bahwa $\rho_q$ memetakan $P$ ke $P$ \emph{(sebab
        kuaternion murni adalah yang $\bar v = -v$)}, bersifat
        linear-$\R$, mengawetkan normanya, dan bahwa $\rho
        \colon q \mapsto \rho_q$ merupakan morfisma grup $S^3
        \to O(3)$.
  \item Hitunglah kernelnya: $\rho_q = \mathrm{id}$ jika dan hanya jika $q$
        komutatif dengan $\mathrm i, \mathrm j, \mathrm k$ jika dan hanya jika
        $q \in \R\cap S^3 = \{\pm1\}$.
  \item Tulislah $q = \cos\frac\theta2 + \sin\frac\theta2\,u$
        dengan $u \in P$ dan $N(u) = 1$ (mengapa ini selalu
        mungkin untuk $q \in S^3$?). Lalu tunjukkan bahwa $\rho_q$ menetapkan
        $u$ dan, pada bidang $u^\perp\cap P$, bekerja sebagai
        rotasi bersudut $\theta$ \emph{(hitunglah
        $\rho_q(w)$ untuk $w \perp u$ dengan memakai $uw = -wu$ bagi
        satuan murni yang ortogonal --- buktikanlah identitas ini dari
        tabel perkaliannya, atau dari $uw + wu =
        -2\langle u, w\rangle$)}.
  \item Simpulkan: bahwa $\operatorname{im}\rho \subseteq SO(3)$
        (sebab setiap $\rho_q$ merupakan rotasi bersumbu dan bersudut seperti
        yang dihitung --- dengan determinan $+1$ lewat \omterm{def:b3:topology:continuity}{kekontinuan} $q
        \mapsto \det\rho_q$ pada $S^3$ yang \omterm{def:b3:topology:connected}{terhubung}, atau
        secara langsung), dan bahwa $\rho$ bersifat \emph{pada} $SO(3)$: sebab setiap
        rotasi $\R^3$ bersumbu (buktikanlah: bahwa matriks ortogonal real
        $3\times3$ dengan $\det = 1$ bernilai eigen $1$ ---
        tinjaulah polinomial karakteristiknya) sehingga berupa suatu $\rho_q$. Ringkasnya:
        \[
        SO(3) \;\cong\; S^3/\{\pm 1\} .
        \]
\end{enumerate}

\medskip
\noindent\textbf{Bagian III --- $SO(3)$ sebagai submanifold;
Rodrigues.}
\begin{enumerate}[resume]
  \item Tunjukkan bahwa $SO(3)$ merupakan submanifold \omterm{def:b3:topology:compact}{kompak}
        berdimensi $3$ di $M_3(\R)$ dengan $T_ISO(3) =$
        matriks antisetangkupnya
        (\cref{exo:b3:submanifolds:5}; sebab syarat determinannya
        memilih sebuah gabungan komponennya).
  \item Untuk matriks antisetangkup $A_u$ yang berkaitan dengan
        $u \in \R^3$ (yakni $A_uv = u\wedge v$, hasil kali
        silangnya), buktikanlah \emph{rumus Rodrigues}:
        \[
        \eu^{\theta A_u} = I + \sin\theta\,A_u + (1 -
        \cos\theta)\,A_u^2 \qquad (\norm u = 1)
        \]
        \emph{(dari $A_u^3 = -A_u$: pisahkanlah deret eksponensialnya
        menurut pangkat $A_u$)}, lalu kenalilah ia sebagai
        rotasi bersumbu $u$ dan bersudut $\theta$. Lalu turunkan bahwa
        $\exp$ memetakan matriks antisetangkupnya \emph{ke seluruh}
        $SO(3)$.
  \item Kaitkanlah kedua parametrisasinya: tunjukkan bahwa $t
        \mapsto \rho_{q(t)}$ dengan $q(t) = \cos\frac t2 +
        \sin\frac t2\,u$ merupakan grup satu parameter berisi
        rotasi yang turunannya di $t = 0$ adalah $A_u$
        --- sehingga eksponensial kuaternion dan matriksnya menuturkan
        cerita yang sama pada laju separuh dan penuh berturut-turut.
\end{enumerate}

\medskip
\noindent\textbf{Bagian IV --- Selimut gandanya, yang terasa.}
\begin{enumerate}[resume]
  \item Tunjukkan bahwa \omterm{def:b3:topology:pathconnected}{lintasan} $q(t) = \cos\frac t2 + \sin\frac
        t2\,\mathrm k$ dengan $t \in \intcc0{2\pi}$ merupakan gelung di
        $SO(3)$ (sebab petanya $\rho_{q(t)}$ kembali ke
        identitasnya) yang angkatan kuaternioniknya \emph{bukan}
        gelung: sebab $q(2\pi) = -q(0)$. Sedangkan melanjutkannya ke $t = 4\pi$
        menutup angkatannya. Lalu terangkanlah dalam sebuah paragraf pendek apa
        yang dikatakannya: bahwa rotasi $2\pi$ tak dapat dibatalkan secara \omterm{def:b3:topology:continuity}{kontinu}
        sedangkan rotasi $4\pi$ dapat (yakni kiat sabuknya),
        sebab gelung $SO(3)$ terdeteksi pada
        selimut gandanya $S^3$.
  \item Turunkan pula buah hasil praktisnya: bahwa komposisi
        rotasi = perkalian kuaternion
        (yakni data senilai $4$ perkalian alih-alih $9$,
        tanpa hanyutan dari keortogonalannya) --- lalu periksalah pada
        komposisi dua seperempat putaran terhadap $\mathrm i$
        dan $\mathrm j$: hitunglah sumbu dan sudut
        hasil kalinya.
\end{enumerate}

\medskip
\noindent\textbf{Bagian V --- Matriks eksplisitnya:
Euler--Rodrigues.} Tulislah $q = a + b\mathrm i + c\mathrm j +
d\mathrm k \in S^3$, sehingga $a^2 + b^2 + c^2 + d^2 = 1$.
\begin{enumerate}[resume]
  \item Hitunglah $\rho_q(\mathrm i)$ selengkapnya dari
        tabel perkaliannya; lalu perolehlah $\rho_q(\mathrm
        j)$ dan $\rho_q(\mathrm k)$ lewat substitusi
        siklik $\mathrm i \to \mathrm j \to \mathrm k
        \to \mathrm i$ dan $(b, c, d) \to (c, d, b)$
        \emph{(benarkanlah: sebab menyiklikkan $\mathrm i, \mathrm j,
        \mathrm k$ diperluas menjadi automorfisma $\mathbb
        H$, karena relasi pendefinisinya setangkup secara
        siklik)}. Lalu simpulkan bahwa matriks $\rho_q$
        pada basis $(\mathrm i, \mathrm j, \mathrm k)$ adalah
        \emph{matriks Euler--Rodrigues}
        \[
        R_q = \begin{pmatrix}
        a^2 + b^2 - c^2 - d^2 & 2(bc - ad) & 2(bd + ac)\\
        2(bc + ad) & a^2 - b^2 + c^2 - d^2 & 2(cd - ab)\\
        2(bd - ac) & 2(cd + ab) & a^2 - b^2 - c^2 + d^2
        \end{pmatrix}.
        \]
  \item (Membaca sebuah rotasi secara mundur) Tunjukkan bahwa
        \[
        \operatorname{tr}R_q = 4a^2 - 1 = 1 + 2\cos\theta,
        \qquad
        \tfrac12\bigl(R_q - R_q^{\mathsf T}\bigr) =
        \sin\theta\,A_u,
        \]
        dalam notasi pertanyaan 5 dan 8. Lalu turunkan sebuah
        algoritma yang memulihkan $\pm q$ dari sebuah matriks rotasi
        $R$: yakni sudutnya dari tracenya; sumbunya dari
        bagian antisetangkupnya bila $0 < \theta < \pi$; dan,
        bila $\theta = \pi$, buktikan lalu pakailah identitas $R
        + I = 2\,uu^{\mathsf T}$.
  \item Nilaikan $R_q$ untuk hasil kali pertanyaan 11 $q =
        \frac12(1 + \mathrm i + \mathrm j + \mathrm k)$: sebuah
        matriks permutasi muncul. Lalu kenali rotasinya
        lalu damaikan dengan sumbu dan sudut yang ditemukan pada
        pertanyaan 11.
\end{enumerate}

\medskip
\noindent\textbf{Bagian VI --- Di dalam $S^3$: $SU(2)$,
\omterm{ex:b3:groups:actions}{kelas konjugasi}, eksponensial.}
\begin{enumerate}[resume]
  \item Tunjukkan bahwa penyajian matriks pertanyaan 1
        (sebutlah $\Phi$) terbatas menjadi isomorfisma grup
        dari $S^3$ ke \emph{grup uniter khusus}
        \[
        SU(2) = \bigl\{U \in M_2(\C) : U^*U = I,\ \det U =
        1\bigr\}
        \]
        \emph{(untuk kesurjektifannya, tuliskanlah persamaan
        $U^{-1} = U^*$ dan $\det U = 1$ bagi sebuah matriks
        kompleks $2\times2$ yang umum)}.
  \item Tunjukkan bahwa bagian realnya merupakan invarian konjugasi
        pada $S^3$ --- yakni $\operatorname{Re}(pq\bar p) =
        \operatorname{Re}q$ untuk setiap $p \in S^3$ --- dan,
        sebaliknya, bahwa dua kuaternion satuan yang berbagian
        real sama bersifat konjugat di $S^3$ \emph{(susutkanlah ke
        pemindahan satu sumbu murni satuan ke sumbu lainnya, yang disediakan Bagian
        II)}. Lalu perikanlah \omterm{ex:b3:groups:actions}{kelas konjugasi}
        $S^3$ secara geometri; terjemahkan ke $SU(2)$ (yakni himpunan
        aras tracenya); lalu proyeksikan lewat $\rho$: bahwa dua
        rotasi bersifat konjugat di $SO(3)$ jika dan hanya jika
        keduanya bersudut $\theta \in \intcc0\pi$ sama.
  \item Definisikan $\exp$ pada $\mathbb H$ lewat deret
        eksponensialnya; periksalah kekonvergenan mutlaknya dengan memakai
        $\abs{pq} = \abs p\,\abs q$ untuk $\abs q =
        \sqrt{N(q)}$. Lalu tunjukkan, untuk $u$ murni satuan dan $\theta
        \in \R$,
        \[
        \exp(\theta u) = \cos\theta + \sin\theta\,u ,
        \]
        lalu turunkan bahwa $\exp$ memetakan hiperbidang $P$ ke
        $S^3$, lalu periksalah bahwa $\rho_{\exp(su)} =
        \eu^{2sA_u}$: yakni gejala setengah sudut pertanyaan
        9 lagi.
  \item Untuk kuaternion murni $v, w$ buktikan aturan hasil kalinya
        $vw = -\langle v, w\rangle + v\wedge w$, sehingga
        identitas \omterm{def:b3:groups:derived}{komutatornya} $vw - wv = 2\,v\wedge w$; lalu buktikan
        pula $[A_v, A_w] = A_{v\wedge w}$ bagi matriks
        pertanyaan 8. Lalu simpulkan bahwa turunan
        $\rho$ di $1$ sepanjang kurva $t \mapsto \exp(tv)$
        adalah isomorfisma linear $v \mapsto 2A_v$ dari $P$
        ke matriks antisetangkupnya, dan bahwa ia
        mengangkut \omterm{def:b3:groups:derived}{komutator} kuaternionnya ke \omterm{def:b3:groups:derived}{komutator}
        matriksnya.
\end{enumerate}

\medskip
\noindent\textbf{Bagian VII --- Struktur globalnya.}
\begin{enumerate}[resume]
  \item (Tak ada seksi \omterm{def:b3:topology:continuity}{kontinu}) Andaikan $s \colon SO(3) \to
        S^3$ \omterm{def:b3:topology:continuity}{kontinu} dengan $\rho \circ s =
        \operatorname{id}$. Lalu untuk gelung $R(t) =
        \rho_{q(t)}$ pada pertanyaan 10, tetapkan $\varepsilon(t) =
        s(R(t))\,q(t)^{-1}$ untuk $t \in \intcc0{2\pi}$. Tunjukkan
        bahwa $\varepsilon$ \omterm{def:b3:topology:continuity}{kontinu} dengan nilai di
        $\{\pm1\}$, lalu turunkan sebuah pertentangan: sehingga tak ada
        pemilihan global yang \omterm{def:b3:topology:continuity}{kontinu} atas kuaternion satuan yang
        mewakili setiap rotasinya.
  \item (Model bolanya) Misalkan $\bar B \subseteq \R^3$ bola
        tertutup berjari-jari $\pi$ dan $E(v) = \eu^{A_v}$,
        dengan $E(0) = I$. Tunjukkan bahwa $E$ memetakan $\bar B$ ke
        $SO(3)$, bersifat injektif pada bola terbukanya, dan pada
        bola batasnya mengenali tepat antipodenya:
        $E(\pi u) = E(-\pi u) = 2uu^{\mathsf T} - I$, tanpa
        kebetulan lainnya. Jadi $SO(3)$ adalah bolanya
        dengan titik batas antipodalnya direkatkan --- yakni
        ruang proyektif $\mathbb{RP}^3$ --- dan sebuah diameternya
        menjadi gelung tak terkontraksikan pertanyaan 10.
  \item Tunjukkan bahwa $\rho_p\rho_q\rho_p^{-1} = \rho_{pq\bar
        p}$; bahwa involusi $SO(3)$ (yakni $R \neq
        I$ dengan $R^2 = I$) tepat berupa setengah putaran
        $\rho_w$ dengan $w$ sebuah kuaternion murni satuan; dan bahwa
        pusat $SO(3)$ bersifat trivial.
  \item Tunjukkan bahwa setiap rotasi merupakan hasil kali dua
        setengah putaran: untuk $q = \cos\frac\theta2 +
        \sin\frac\theta2\,u$, pilihlah sebuah murni satuan $w \perp
        u$, periksalah bahwa $w' = qw$ kembali merupakan kuaternion
        murni satuan, lalu periksalah $\rho_q =
        \rho_{w'}\rho_w$. Lalu di manakah kedua sumbunya terletak, dan
        bersudut berapa keduanya?
  \item Simpulkan ringkasan \omterm{def:b3:topology:topology}{topologinya}: bahwa $SO(3)$ bersifat
        \omterm{def:b3:topology:compact}{kompak} dan \omterm{def:b3:topology:pathconnected}{terhubung lintasan} (berikan dua bukti:
        yakni peta \omterm{def:b3:topology:continuity}{kontinu} $S^3$ di bawah $\rho$; dan peta
        $\exp$), sedangkan $O(3)$ tepat berkomponen \omterm{def:b3:topology:connected}{terhubung}
        dua, yang masing-masingnya \omterm{def:b3:topology:continuity}{homeomorfik} dengan $SO(3)$.

  \item (Sebuah komposisi, tiga cara) Misalkan $R_1$ rotasi
        sebesar $\frac\pi2$ terhadap sumbu-$z$ dan $R_2$
        rotasi sebesar $\frac\pi2$ terhadap sumbu-$x$. Hitunglah
        sumbu dan sudut $R_2R_1$: (i) dengan mengalikan
        kedua matriks $3\times3$-nya lalu memakai
        trace dan bagian antisetangkupnya (Bagian V); (ii) dengan
        mengalikan kuaternion satuan yang bersesuaian
        $q_2q_1$. Lalu periksalah bahwa kedua jawabannya bersesuaian: sudutnya
        $\frac{2\pi}3$ dan sumbunya $\frac1{\sqrt3}(1, -1, 1)$.
  \item (\omterm{ex:b3:conformal:mobius}{Transformasi Cayley}) Untuk $K$ yang antisetangkup, tunjukkan
        bahwa $I + K$ terbalikkan dan
        \[
        C(K) = (I - K)(I + K)^{-1} \in SO(n),
        \]
        dengan $-1$ yang tak pernah menjadi nilai eigen $C(K)$; lalu tunjukkan bahwa
        $K \mapsto C(K)$ merupakan bijeksi dari matriks antisetangkup
        ke $\{R \in SO(n) : -1 \notin
        \operatorname{Sp}R\}$, dengan balikan $R \mapsto (I -
        R)(I + R)^{-1}$. (Yakni sebuah peta rasional $SO(n)$, yang
        mendampingi $\exp$ transenden pada
        \cref{exo:b3:submanifolds:11}.)

\end{enumerate}
\end{problem}
